\section{Introducción}
\label{sec:introduccion}

\subsection{Multiestabilidad y localización de atractores}

Un mismo sistema dinámico puede presentar respuestas muy diferentes sin
cambiar sus parámetros. Un conjunto de condiciones iniciales puede conducir
a un equilibrio, mientras que otro produce una oscilación sostenida o una
respuesta caótica. Esta coexistencia, conocida como \emph{multiestabilidad},
hace que la posición inicial sea tan relevante como las ecuaciones para
comprender el comportamiento observado. En electrónica, interesa tanto por
la posibilidad de seleccionar regímenes de operación como por la aparición
de oscilaciones que un análisis alrededor del punto de funcionamiento puede
pasar por alto \cite{DudkowskiEtAl2016,WangKuznetsovChen2021Multistability}.

La noción de atractor permite describir las respuestas hacia las que evoluciona
un conjunto de estados iniciales. Su cuenca de atracción reúne los datos que
alcanzan esa respuesta a largo plazo. Conocer un atractor y conocer su cuenca
son problemas relacionados, pero diferentes: una trayectoria puede mostrar
con claridad una oscilación sin revelar desde dónde es accesible. Los textos
de sistemas dinámicos desarrollan esta relación entre estabilidad local,
comportamiento asintótico y organización del espacio de fases
\cite{HirschSmaleDevaney2013,Strogatz2015Nonlinear}.

La distinción entre atractores autoexcitados y ocultos surge precisamente de
la relación de sus cuencas con los equilibrios. Un atractor \(\mathcal A\) es
\emph{autoexcitado} cuando su cuenca intersecta toda vecindad de, al menos,
un equilibrio; es \emph{oculto} cuando, para cada equilibrio, existe una
vecindad que no intersecta esa cuenca
\cite{LeonovKuznetsov2013,DudkowskiEtAl2016}. En el primer caso, comenzar
suficientemente cerca de un equilibrio adecuado puede proporcionar una
trayectoria que llegue al atractor. En el segundo, la búsqueda debe alcanzar
otra parte del espacio de estados. El término oculto se refiere a esta
dificultad de acceso desde los equilibrios: puede describir tanto una
oscilación regular como un atractor caótico.

Esta diferencia cambia la manera de estudiar un modelo. Si todas las
simulaciones se inician cerca de sus equilibrios, algunas respuestas pueden
quedar sin observar, incluso cuando el modelo sea de dimensión baja. Por
otro lado, encontrar una trayectoria interesante mediante prueba y error
deja abierta una pregunta práctica: cómo volver a encontrarla al variar
los parámetros, cambiar el orden de derivación o aplicar la búsqueda a
otro sistema. La localización analítico--numérica aborda esa pregunta al
construir condiciones iniciales a partir de la estructura de las ecuaciones
\cite{LeonovKuznetsovVagaitsev2011,KuznetsovEtAl2017}.

La motivación de esta investigación parte de ese problema. Buscamos que la
obtención de una respuesta dinámica pueda explicarse desde el modelo: qué
información permite proponer una condición inicial, cómo se transporta esa
condición hasta el sistema de interés y cómo se identifica el destino
alcanzado. La selección de la condición inicial deja así de ser un dato
aislado y se convierte en una parte del método que puede estudiarse y
compararse.

\subsection{Orden fraccionario, memoria y aplicación digital}

Los sistemas de orden fraccionario amplían este problema al incorporar
operadores cuya acción depende de un intervalo de la evolución. En esta
tesis trabajamos principalmente con la derivada de Caputo, adecuada para
formular problemas de valor inicial con condiciones puntuales habituales.
Para \(0<q<1\), su acción sobre una función suficientemente regular es
\[
{}^{\mathrm C}D_{t_0}^{q}X(t)
=\frac{1}{\Gamma(1-q)}
\int_{t_0}^{t}(t-\tau)^{-q}X'(\tau)\,\mathrm d\tau.
\]
Aquí \(q\) es el orden, \(t_0\) el instante inicial y \(\Gamma\) la función
gamma. La integral muestra cómo intervienen los valores anteriores de la
evolución: el peso depende de la distancia temporal \(t-\tau\). El efecto
de memoria pertenece, por tanto, a la ecuación que se pretende resolver
\cite{Podlubny1999,Diethelm2010AnalysisFDE}.

Esta propiedad ofrece una forma de estudiar respuestas que dependen de su
evolución previa, pero también obliga a precisar qué ocurre durante una
búsqueda. En una ecuación diferencial ordinaria, el estado en un instante
intermedio permite continuar la solución. En un problema de Caputo con
terminal inferior fijo, continuar desde ese instante requiere conservar la
contribución de la historia anterior. Si se toma únicamente el último
vector de estado como un nuevo dato inicial, se define otro problema.
La diferencia resulta especialmente relevante al encadenar etapas de
continuación o comparar algoritmos de localización
\cite{Diethelm2020_afcee,DoanKloeden2021Semidynamical}.

El orden también influye en la estabilidad. En el caso lineal conmensurable,
la condición depende del ángulo de los valores propios respecto del eje real
positivo, mediante el criterio de Matignon. La extensión de ese análisis a
un sistema no lineal requiere además las hipótesis correspondientes de
linealización \cite{Matignon1996,CongDoanSiegmundTuan2016}. Por ello,
modificar \(q\) no equivale simplemente a recorrer más despacio las
trayectorias del sistema entero: puede cambiar la estabilidad y la
coexistencia de respuestas.

La aplicación que da contexto a la tesis es el cifrado de imágenes mediante
secuencias producidas por sistemas fraccionarios multiestables. El interés
consiste en estudiar cómo seleccionar y reproducir un régimen dinámico y
cómo convertir su evolución en una secuencia digital. La realización de
sistemas fraccionarios en DSP y de generadores pseudoaleatorios en ESP32 y
FPGA ofrece antecedentes concretos de esa transición entre modelo,
algoritmo e implementación
\cite{Wu2023ArctanChua,Fraga2026_54625}.

En esa transición, las propiedades del modelo continuo deben relacionarse
con las de su representación digital. La precisión aritmética, la
aproximación de memoria, el muestreo y la conversión a enteros intervienen
en la secuencia obtenida. La multiestabilidad añade la necesidad de asegurar
que la inicialización selecciona la respuesta buscada y que esa selección
puede repetirse. Así, el estudio de los atractores y sus cuencas proporciona
una base para elegir el generador y sus condiciones de operación, mientras
que la evaluación del cifrado examina la transformación de la imagen y su
resistencia a ataques \cite{Alvarez2006_3e24c,Lima2025_99c51}.

La relación entre estas tareas orienta el trabajo: antes de utilizar una
secuencia, necesitamos comprender qué dinámica la produce, cómo se obtiene
y qué decisiones numéricas afectan su reproducción. Después podemos
examinar su utilidad en una aplicación digital concreta. Esta secuencia de
preguntas conecta el problema matemático de localización con la finalidad
electrónica de la investigación.

\section{Antecedentes y estado del arte}
\label{sec:antecedentes}

El circuito de Chua ocupa un lugar importante en el estudio experimental y
matemático del caos por la relación directa entre una realización electrónica
de pocos elementos y un modelo no lineal de baja dimensión. Matsumoto mostró
en 1984 un atractor caótico de este circuito
\cite{Matsumoto1984Chua}. Años después, la búsqueda de respuestas que no se
obtienen al iniciar cerca de sus equilibrios abrió otra línea de estudio
dentro de la misma familia. Leonov, Kuznetsov y Vagaitsev localizaron
atractores ocultos de Chua en 2011, y el análisis posterior de una
característica suave mostró que el fenómeno también aparece fuera de los
modelos con característica lineal por tramos
\cite{LeonovKuznetsovVagaitsev2011,Leonov2013HiddenAttractors}.

El interés del problema va más allá de un circuito particular. La revisión
de Leonov y Kuznetsov relaciona las oscilaciones ocultas con problemas
clásicos de estabilidad absoluta, entre ellos las conjeturas de Aizerman y
Kalman \cite{LeonovKuznetsov2013}. Estos problemas muestran por qué la
estabilidad de una familia de aproximaciones lineales no basta, en general,
para describir toda la dinámica del sistema no lineal. Las dos conjeturas
imponen condiciones diferentes sobre la característica de realimentación;
comparten la dificultad de pasar de una información local o sectorial a
una conclusión global. Un ejemplo posterior de la familia Kalman--Fitts
presenta coexistencia de un equilibrio localmente estable y ciclos ocultos
\cite{kalman2019}.

Desde la perspectiva de ingeniería, esa coexistencia significa que un
funcionamiento aparentemente estable puede depender de la inicialización
o de la magnitud de una perturbación. En los lazos de amarre de fase,
Leonov y colaboradores documentan cómo ciertas simulaciones y modelos
simplificados pueden omitir oscilaciones relevantes para interpretar el
funcionamiento del dispositivo \cite{pll2015}. La localización de
atractores ocultos responde, por tanto, a una necesidad de análisis de
regímenes coexistentes, además del interés por construir nuevos ejemplos.

\subsection{Leonov y Kuznetsov: método en orden entero}

Leonov, Kuznetsov y colaboradores formularon un procedimiento
analítico--numérico para localizar atractores ocultos en sistemas tipo Lur'e
\cite{LeonovKuznetsov2013,KuznetsovEtAl2017}. La separación en un bloque lineal y una no linealidad escalar permite
calcular la transferencia. Al proyectar la respuesta no lineal sobre su
componente fundamental obtenemos una condición de balance armónico. La
frecuencia y la ganancia resultantes definen un sistema lineal auxiliar y
un modo oscilatorio. A partir de ese modo construimos la condición inicial
y la transportamos por continuación hasta el campo original.

El balance armónico combina la transferencia y la función descriptiva para
obtener una frecuencia, una ganancia equivalente y una amplitud. Estos
valores permiten construir la semilla; el estudio posterior de su cuenca
establece si la respuesta localizada es accesible desde las vecindades
de los equilibrios.

La ventaja de esta construcción es que aprovecha la estructura de
realimentación del modelo. En lugar de integrar una gran cantidad de
condiciones iniciales sin orientación, utiliza la respuesta frecuencial
de la parte lineal y una aproximación de la no linealidad para proponer
una oscilación inicial. Su aplicabilidad depende de que exista una
realización adecuada, de que el cierre armónico admita una solución y
de que la continuación alcance el sistema original. Cuando una de esas
condiciones falla, la dificultad aporta información sobre la estructura
del caso y motiva una ruta de localización diferente.

\subsection{Danca: Chua fraccionario no suave}

Danca estudió atractores ocultos en varios sistemas fraccionarios, incluido un
sistema de Chua con saturación \cite{Danca2017}. El artículo usa derivadas
de Caputo de orden conmensurable, clasifica los equilibrios mediante Matignon e integra
el problema con un predictor--corrector Adams--Bashforth--Moulton. Para el
ejemplo de Chua se reportan
\[
q=0.9998,\qquad
(\alpha,\beta,\gamma)=(8.4562,12.0732,0.0052),\qquad
(m_0,m_1)=(-0.1768,-1.1468).
\]
La verificación descrita en la referencia usa 200 trayectorias alrededor de un
equilibrio externo con radio \(\delta=0.01\) y emplea la simetría para el
equilibrio opuesto.

La reconstrucción conserva la ecuación, los
parámetros, el orden y el integrador de la referencia. El artículo no
proporciona la condición inicial exacta de la figura del atractor ni los
valores intermedios de la construcción armónica
\((\omega_0,k,A_0)\). Por ello, la reconstrucción utiliza la semilla obtenida
al aplicar el cierre armónico entero a los parámetros publicados.
La ruta propuesta usa los mismos parámetros
nominales, pero calcula una nueva semilla sesgada y aplica una continuación
compatible con la memoria de Caputo.

\subsection{Wu et al.: Chua fraccionario con arctangente}

Wu et al. propusieron un sistema de Chua fraccionario con no linealidad
arctangente \cite{Wu2023ArctanChua}. La referencia emplea una función
descriptiva, una recurrencia local de descomposición de Adomian y las
condiciones iniciales opuestas
\[
X_{0,+}=(13.8,\;0.7093,\;-19.8768)^{\T},
\qquad X_{0,-}=-X_{0,+}.
\]
Los parámetros publicados se expresan mediante
\[
h(x)=mx+(n-m)\arctan(x),\qquad
m=0.4,\qquad n=-1.1585,
\]
de modo que, en la notación empleada aquí,
\[
a_1=0.4,\qquad a_2=n-m=-1.5585,\qquad \rho=1,\qquad q=0.99.
\]

La reconstrucción bibliográfica reproduce la recurrencia ADM local y las dos
condiciones iniciales publicadas. La ruta propuesta para el sistema con
arctangente es independiente: explora un dominio acotado de parámetros en
\(q=1\), selecciona un caso mediante reglas reproducibles y refina la semilla
con problemas de Caputo de memoria completa.

Los antecedentes de Danca y Wu permiten formular una comparación con
modelos concretos y parámetros publicados. También muestran que dos
trabajos descritos como fraccionarios pueden construir sus trayectorias
mediante procedimientos distintos. Por ello, la comparación de resultados
debe considerar conjuntamente la no linealidad, el operador, la
inicialización y la forma en que el método numérico trata la memoria.

\subsection{Memoria, estabilidad y periodicidad en el marco fraccionario}

La formulación de Caputo puede expresarse como una ecuación integral de
Volterra. Para el problema
\[
{}^{\mathrm C}D_{t_0}^{q}X(t)=F(X(t)),\qquad X(t_0)=X_0,
\qquad 0<q<1,
\]
donde \(F\) es el campo vectorial y \(X_0\) el estado inicial, la expresión
integral equivalente es
\[
X(t)=X_0+\frac{1}{\Gamma(q)}
\int_{t_0}^{t}(t-\tau)^{q-1}F(X(\tau))\,\mathrm d\tau.
\]
El valor de \(X(t)\) depende de las contribuciones del campo a lo largo de
todo el intervalo. Esta formulación es la base de numerosos métodos
numéricos y permite identificar qué información se conserva cuando se
continúa una solución \cite{Diethelm2010AnalysisFDE,DiethelmFordFreed2002}.

El predictor--corrector de Adams--Bashforth--Moulton de Diethelm, Ford y
Freed proporciona una referencia para integrar estos problemas
\cite{DiethelmFordFreed2002}. Las aproximaciones de memoria corta y los
métodos rápidos buscan reducir el costo de las sumas de historia. Su
pertinencia depende de cómo aproximan el núcleo y de sus controles de error:
truncar una historia, comprimirla y comenzar un nuevo problema desde el
último estado son operaciones diferentes
\cite{Deng2007,KhristenkoWohlmuth2023Rational,JiangZhangZhangZhang2017FastCaputo}.
Diethelm, Garrappa y Stynes explican además que la regularidad de la solución
cerca del instante inicial afecta el orden de convergencia, y que una
recurrencia local no reproduce automáticamente un operador no local
\cite{Diethelm2020_afcee}.

El estudio de atractores exige precisar también el espacio en el que se
considera la evolución. Los trabajos de Doan y Kloeden construyen marcos
semidinámicos para ecuaciones autónomas de Caputo y desarrollan el análisis
de sus atractores \cite{DoanKloeden2021Semidynamical,DoanKloeden2024Attractors}.
Estos antecedentes permiten relacionar el estado observable con la
información de memoria necesaria para describir una continuación. En la
experimentación, fijar una inicialización común hace posible comparar las
cuencas obtenidas para una misma clase de datos iniciales.

La periodicidad constituye otro punto de encuentro entre la teoría y la
localización armónica. Los resultados de Tavazoei y Haeri, y de Kaslik y
Sivasundaram, establecen restricciones para soluciones periódicas exactas
en formulaciones fraccionarias causales bajo sus respectivas hipótesis
\cite{TavazoeiHaeri2009,KaslikSivasundaram2012}. Una oscilación que parece
cerrarse después de un transitorio no identifica por sí sola la misma
noción de periodicidad que se emplea en una ecuación diferencial ordinaria.
La diferencia importa cuando se utiliza el balance armónico para construir
una condición inicial.

Los operadores de Weyl permiten formular el análisis de señales periódicas
con otra historia temporal. Vigué y colaboradores estudian la continuación
de soluciones periódicas en ese marco; Haacker y colaboradores desarrollan
un análisis de estabilidad de tipo Hill para la formulación de
Liouville--Weyl \cite{Vigue2019,Haacker2026}. En un oscilador lineal forzado,
Duan, Lan y Li identifican en la respuesta de Caputo una componente
periódica de Weyl acompañada de términos transitorios, incluida una cola
algebraica \cite{DuanLanLi2022WeylCaputo}. Este ejemplo explica por qué la
aproximación armónica puede ser informativa sin identificarla con toda la
evolución causal del sistema no lineal. Los estudios sobre derivadas de
funciones periódicas y casi periódicas completan esta distinción entre
operador, dominio temporal y respuesta
\cite{OrtigueiraMachadoTrujillo2017Periodic,ZhuLi2026PeriodicAlmostPeriodic}.

\subsection{Métodos geométricos y estudio de las cuencas}

La función descriptiva es una de varias formas de obtener información para
la búsqueda. Los puntos perpetuos, introducidos por Prasad, se definen
mediante una condición de aceleración nula del campo y velocidad no nula.
En un sistema entero \(\dot X=F(X)\), esa condición se escribe
\[
DF(p)F(p)=0,\qquad F(p)\neq0.
\]
Su interés consiste en proponer estados especiales del campo que pueden
emplearse como condiciones iniciales
\cite{Prasad2015Perpetual,DudkowskiPrasadKapitaniak2015,DudkowskiPrasadKapitaniak2017}.
La pregunta relevante es qué destinos alcanzan esas semillas en cada
modelo. Los estudios de suficiencia y los contraejemplos muestran que
su utilidad depende del sistema, por lo que resulta natural compararlos
con otras construcciones bajo las mismas condiciones de integración
\cite{NazarimehrSaediJafari2017}.

Las curvas conectantes y las superficies críticas extraen otra información
de las derivadas del campo. Guan y Xie utilizan curvas conectantes para
orientar la localización, mientras que Godara y colaboradores estudian
estructuras topológicas asociadas con sistemas multiestables
\cite{GuanXie2021Connecting,GodaraEtAl2018CriticalSurfaces}.
Estas construcciones son especialmente interesantes cuando una
representación de realimentación escalar no resulta suficiente o cuando el
cierre armónico directo no tiene solución. Los ejemplos de
Rabinovich--Fabrikant amplían además el estudio de dinámicas y atractores
ocultos fuera de la familia de Chua \cite{Danca2017HiddenFO}.

Una vez localizada una respuesta, la cuenca permite estudiar su accesibilidad
y la coexistencia con otros destinos. Sprott y Xiong proponen clasificar y
cuantificar cuencas considerando su extensión y la distribución de
condiciones iniciales \cite{SprottXiong2015Basins}. La entropía de cuenca
de Daza y colaboradores relaciona la incertidumbre del destino con la
resolución y la coexistencia de respuestas
\cite{DazaEtAl2016BasinEntropy}. El cálculo computacional de cuencas y la
identificación de atractores mediante recurrencias ofrecen herramientas
adicionales para explorar estos problemas \cite{datseris2022}.

Estas contribuciones sitúan el análisis de cuencas como una parte de la
localización. Un corte plano facilita observar la distribución de destinos;
los sondeos alrededor de los equilibrios responden a la clasificación de
ocultedad; una exploración más amplia describe regiones de accesibilidad.
Las tres observaciones aportan información complementaria. La resolución,
el dominio y el horizonte deben acompañar la interpretación de cada
figura, especialmente si se estudian fronteras entremezcladas o propiedades
como Wada \cite{DazaEtAl2015Wada}.

\subsection{Diagnósticos dinámicos y comparación numérica}

La caracterización de una trayectoria combina su geometría y su evolución
temporal. Proyecciones del espacio de estados, secciones de Poincaré,
espectros y estimaciones de exponentes de Lyapunov responden a preguntas
distintas sobre recurrencia, frecuencias y sensibilidad. En orden entero,
los métodos variacionales de Benettin y colaboradores proporcionan un
fundamento clásico para calcular exponentes
\cite{Benettin1980a,Benettin1980b}. Para los sistemas fraccionarios se han
propuesto adaptaciones variacionales, mapas de Poincaré y otros
diagnósticos \cite{DancaKuznetsov2018Lyapunov,fischer2020,ClementeLopezEtAl2022Poincare}.

En Caputo, el tratamiento de memoria interviene también en el problema
variacional. Danca y Fečkan identifican un esquema impulsivo con terminal
inferior cambiante en el código que analizan; una propuesta posterior
incorpora QR y el integrador LIL manteniendo esa interpretación
\cite{DancaFeckan2024MemoryPrinciple,Danca2026ImprovedLyapunov}.
Por ello, al comparar estimaciones es necesario relacionar cada número con
su ecuación variacional, su intervalo de observación y su política de
renormalización. Esta precisión permite interpretar discrepancias entre
algoritmos sin atribuirlas de inmediato a un cambio del sistema físico.

La comparación numérica tiene una función científica dentro de la
propuesta. Reducir el paso y ampliar el horizonte permite examinar la
persistencia de la respuesta; utilizar otra implementación ayuda a
detectar decisiones que dependen del algoritmo. La revisión de buenas
prácticas de cálculo fraccionario proporciona el contexto para estos
contrastes \cite{Diethelm2020_afcee}. El propósito es establecer qué
características de la dinámica se reproducen bajo condiciones comparables
y qué diferencias merecen un análisis adicional.

\subsection{Multiestabilidad, generación de secuencias y cifrado de imágenes}

Las aplicaciones de los atractores ocultos se han desarrollado junto con
su estudio teórico y su realización electrónica. Las monografías sobre
sistemas con atractores autoexcitados y ocultos reúnen ejemplos y
construcciones en distintos modelos, incluidas implementaciones analógicas
y digitales \cite{PhamEtAl2018HiddenAttractors,PhamVolosKapitaniak2017Systems}.
El trabajo experimental sobre Chua vincula directamente la teoría de
localización con la observación de respuestas del circuito
\cite{KuznetsovEtAl2023ChuaExperiments}. Esta conexión resulta relevante
para una tesis en electrónica: un modelo localizado numéricamente debe
poder relacionarse con una realización y unas condiciones de operación.

En el ámbito fraccionario, Wu y colaboradores incluyen una implementación
DSP del sistema con arctangente \cite{Wu2023ArctanChua}. De la Fraga y
Tlelo-Cuautle estudian la implementación de un generador pseudoaleatorio
basado en Lorenz fraccionario en ESP32 y FPGA
\cite{Fraga2026_54625}. Son antecedentes de tareas distintas: localizar
y representar una respuesta dinámica, por una parte, y producir una
secuencia digital mediante una arquitectura determinada, por otra.
Su relación motiva estudiar conjuntamente el modelo, el integrador y la
extracción de la secuencia, en lugar de elegirlos de manera independiente.

La precisión finita introduce una diferencia esencial. Si una realización
digital es determinista y autónoma, y su estado completo, incluida la
memoria del algoritmo, admite sólo un número finito de valores, alguna
configuración debe repetirse si la ejecución se prolonga suficientemente.
A partir de esa repetición, la evolución vuelve a recorrer los mismos
estados. Por ello, la ausencia de periodicidad del modelo continuo no
garantiza por sí sola las propiedades de la secuencia implementada.
Álvarez y Li sitúan la descripción de la implementación, el manejo de claves
y el análisis de seguridad entre los requisitos básicos de un criptosistema
basado en caos \cite{Alvarez2006_3e24c}. Lima y colaboradores estudian una
arquitectura que aprovecha diferencias de precisión finita entre dos
realizaciones digitales de Chua para el cifrado de imágenes
\cite{Lima2025_99c51}; ese resultado muestra que el efecto aritmético puede
formar parte de una construcción explícita y debe analizarse como tal.

La multiestabilidad ofrece varios regímenes posibles con los mismos
parámetros, lo que hace interesante la selección de condiciones iniciales
para producir secuencias. A la vez, exige identificar qué regiones llevan
a equilibrios, oscilaciones regulares o respuestas caóticas. El tamaño
nominal de un conjunto de parámetros y condiciones iniciales no equivale
automáticamente al número de claves distinguibles y útiles de una
implementación \cite{Alvarez2006_3e24c}. El estudio de cuencas permite
relacionar esa selección con el comportamiento del generador.

La evaluación estadística de una secuencia y la seguridad del cifrado
aportan evidencias diferentes. Zhao, Shi y Ding presentan un ataque de
texto claro elegido contra un algoritmo específico que combina caos y
codificación de ADN, aun cuando sus pruebas estadísticas sugerían buenas
propiedades \cite{Zhao2025_98e42}. Este antecedente justifica incluir
un análisis de la estructura de transformación y del escenario de ataque
cuando se estudia una aplicación propia. La tesis conecta así localización,
generación digital y cifrado mediante preguntas que pueden evaluarse
por separado y relacionarse después.

\subsection{Síntesis de los métodos y problema pendiente}

La revisión reciente de Guan y Xie reúne métodos analíticos, numéricos y
geométricos de localización \cite{GuanXie2025}. El campo dispone, por tanto,
de varias rutas de búsqueda, tanto para sistemas clásicos como para
familias fraccionarias. La cuestión que interesa aquí es cómo relacionar
esas rutas con una formulación causal y con una comparación que conserve
el mismo problema de evolución. Los tres antecedentes principales de
Chua proporcionan un punto de partida concreto para esa relación:

\begin{table}[H]
\centering
\caption{Papel de los trabajos de referencia en la metodología.}
\label{tab:papel_referencias}
\begin{tabularx}{\linewidth}{L{0.22\linewidth}Y Y}
\toprule
Referencia & Elemento recuperado & Función metodológica\\
\midrule
Leonov--Kuznetsov
& Función descriptiva, construcción modal y continuación en \(q=1\)
& Caso límite y fundamento de la ruta de localización\\
Danca
& Chua no suave, Caputo, Matignon y ABM
& Reconstrucción bibliográfica y sistema base de la ruta sesgada\\
Wu et al.
& Chua con arctangente, condiciones iniciales y ADM local
& Reconstrucción bibliográfica y comparación con la ruta independiente
\texttt{c590}\\
\bottomrule
\end{tabularx}
\end{table}

\subsection{Análisis matemático de los artículos primarios}

\begin{longtable}{@{}L{0.16\textwidth}L{0.22\textwidth}L{0.28\textwidth}L{0.26\textwidth}@{}}
\caption{Operaciones matemáticas asociadas con el estado del arte.}
\label{tab:articulos-operaciones}\\
\toprule
Núcleo & Artículos & Formalismo reconstruido & Alcance de la interpretación \\
\midrule
\endfirsthead
\toprule
Núcleo & Artículos & Formalismo reconstruido & Alcance de la interpretación \\
\midrule
\endhead
Atractor oculto y localización
& \cite{LeonovKuznetsovVagaitsev2011,LeonovKuznetsov2013,KuznetsovEtAl2017}
& Realización de Lur'e; resolvente; función descriptiva; identidad modal;
homotopía desde un sistema auxiliar.
& La construcción produce una semilla. Atracción, caos y ocultedad requieren
integración y prueba de cuenca separadas. \\
Caputo y atractores ocultos
& \cite{Danca2017}
& PVI de Caputo; equilibrios; jacobiano; criterio de Matignon; PECE; sondeos en
vecindades.
& El espectro es local y los sondeos son finitos. La memoria y el terminal
inferior forman parte del problema. \\
Operador y método numérico
& \cite{Matignon1996,DiethelmFordFreed2002}
& Transformada de Laplace; Mittag--Leffler; forma de Volterra; pesos predictor
y corrector; orden de convergencia.
& Matignon no clasifica cuencas. PECE aproxima un PVI declarado; no corrige una
historia reiniciada o indefinida. \\
Función descriptiva fraccionaria
& \cite{BarbosaMachado2002BacklashImpact,Machado2015FDF}
& Coeficiente fundamental complejo; dependencia de frecuencia; ramas de
\(N^\mu\); cierre auxiliar y cierre físico.
& Es un predictor armónico. La potencia de Machado puede diversificar semillas,
pero no sustituye la proyección de Fourier de la no linealidad física. \\
Puntos perpetuos
& \cite{Prasad2015Perpetual,DudkowskiPrasadKapitaniak2017,NazarimehrSaediJafari2017}
& \(Df(p)f(p)=0\), exclusión de equilibrios, singularidad de \(Df\),
covariancia afín y contraejemplos.
& Un PP es una semilla heurística, no un localizador universal. FPP--A y FPP--H
son formulaciones propuestas para Caputo y deben conservar el formalismo de
memoria. \\
Geometría del campo
& \cite{GuanXie2021Connecting,GodaraEtAl2018CriticalSurfaces}
& Menores, superficies \(\det Df=0\), campos derivados, curvas conectantes y
transversalidad.
& Estos objetos organizan semillas y regiones; no prueban por sí solos
invariancia, atracción ni una frontera global. \\
Topología de cuencas
& \cite{DazaEtAl2016BasinEntropy,DazaEtAl2015Wada}
& Partición del dominio, entropía de cuenca, refinamiento y criterio de Wada.
& Una imagen entremezclada es evidencia finita. La afirmación topológica exige
convergencia o caracterización del conjunto de frontera. \\
Dinámica con memoria
& \cite{DoanKloeden2024Attractors}
& Semidinámica en un espacio de historias, invariancia funcional, atracción y
proyección al estado observable.
& La geometría puntual de una ODE no se traslada a Caputo mediante una regla de
cadena local ni mediante reinicios sucesivos. \\
\bottomrule
\end{longtable}

La comparación permite identificar un problema metodológico concreto.
Una semilla puede derivarse con precisión y conducir a un destino distinto
del esperado; una aproximación armónica puede ser útil para iniciar una
búsqueda sin describir toda la evolución de Caputo; una respuesta puede
persistir en un integrador y cambiar al modificar la discretización. Estas
posibilidades forman parte del problema de localización y deben examinarse
en una misma secuencia de análisis. De otro modo, es difícil saber si dos
procedimientos difieren por la información que utilizan para buscar o
porque integran problemas distintos.

La propuesta de esta tesis consiste en conectar la generación estructurada
de semillas con la evolución causal, la identificación de destinos y el
estudio de cuencas. El balance armónico, la continuación y la geometría
del campo aportan información complementaria. El objetivo es determinar
cuándo resulta útil cada información, cómo combinarlas y cómo comparar
su desempeño con una exploración directa del mismo sistema. Este
planteamiento conserva la función de los métodos publicados y permite
estudiar las modificaciones propuestas mediante casos y contrastes
específicos.

\subsection{Problema y preguntas de investigación}
\label{sec:cap1-problema}

El problema central es construir un procedimiento que permita localizar
y caracterizar atractores ocultos en sistemas multiestables de orden
fraccionario, relacionando la estructura matemática del modelo con las
condiciones iniciales que alcanzan sus respuestas. El interés no se limita
a obtener una trayectoria: necesitamos explicar la ruta que lleva a ella,
examinar su estabilidad frente a decisiones numéricas y establecer cómo
se relaciona su cuenca con los equilibrios.

La pregunta que guía la investigación es: \emph{¿cómo combinar información
de realimentación, aproximaciones armónicas y geometría del campo para
localizar atractores ocultos en sistemas fraccionarios, conservando una
formulación causal de memoria y obteniendo resultados que puedan
reproducirse y compararse?} La aplicación al cifrado añade la pregunta de
cómo utilizar las respuestas seleccionadas para producir secuencias
digitales y evaluar una transformación reversible de imágenes.

Esta pregunta se desarrolla en cuatro problemas relacionados:
\begin{enumerate}
\item ¿Qué propiedades del modelo permiten construir una semilla y en qué
      casos una obstrucción del cierre armónico exige utilizar otra ruta?
\item ¿Cómo transportar o refinar esa semilla sin confundir una
      continuación con memoria acumulada y una sucesión de problemas
      iniciales independientes?
\item ¿Qué respuestas se localizan, con qué costo y con qué relación
      respecto de las vecindades de los equilibrios, al comparar métodos
      bajo un presupuesto y un protocolo comunes?
\item ¿Cómo influyen la selección del régimen, la discretización y la
      precisión finita en la generación de secuencias y en su utilización
      para el cifrado de imágenes?
\end{enumerate}

\subsection{Hipótesis de trabajo y contribución metodológica}
\label{sec:cap1-hipotesis}

La hipótesis de trabajo es que la información armónica y la geometría del
campo permiten orientar de manera complementaria la selección de
condiciones iniciales, de modo que una búsqueda estructurada puede
localizar respuestas que no se observan al iniciar cerca de los equilibrios.
En sistemas de Caputo, esa información debe integrarse con un tratamiento
explícito de la historia para que la comparación corresponda al mismo
problema de evolución. La utilidad de la combinación se contrasta mediante
los destinos alcanzados, la persistencia de las respuestas y el costo
computacional, frente a rutas alternativas con condiciones comparables.

La aportación se centra en la relación entre esas etapas. Para el Chua
no suave estudiamos una función descriptiva sesgada y una continuación
con historia acumulada. Para el Chua con arctangente consideramos una
exploración inicial en orden entero y un refinamiento posterior mediante
problemas independientes de Caputo. Los casos Kalman--Fitts, MAVPD y PLL
permiten comparar la búsqueda armónica con mapas de conmutación,
homotopías y semillas geométricas. Las extensiones a otras familias
fraccionarias permiten estudiar hasta dónde se conserva esa utilidad
cuando cambia la estructura del campo.

El alcance general se apoya en distinguir qué parte de la construcción
depende de una realimentación escalar y cuál puede formularse mediante
varios canales o mediante propiedades geométricas. La comparación incluye
casos donde una ruta funciona y casos donde no alcanza el destino
buscado. Esta diversidad permite identificar condiciones de aplicación y
explicar por qué una misma fórmula de inicialización puede producir
resultados diferentes en dos modelos.

\subsection{Objetivo general y objetivos específicos}
\label{sec:cap1-objetivos}

El objetivo general es desarrollar y evaluar una metodología matemática
y computacional para localizar y caracterizar atractores ocultos en
sistemas multiestables de orden fraccionario, utilizando el orden entero
como referencia y relacionando la estructura del modelo, la memoria de
Caputo y las cuencas de atracción. A partir de los sistemas seleccionados,
se estudia su utilización para generar secuencias digitales y para una
aplicación de cifrado de imágenes.

Los objetivos específicos son:
\begin{enumerate}
\item Formular los modelos de estudio, calcular sus equilibrios y
      jacobianos, y establecer la clasificación de estabilidad local
      correspondiente al orden de derivación y a las hipótesis de cada caso.
\item Construir y analizar condiciones iniciales mediante realizaciones
      de realimentación, función descriptiva, balance armónico y métodos
      geométricos complementarios, identificando sus condiciones de
      admisibilidad y sus posibles obstrucciones.
\item Desarrollar las rutas de continuación y refinamiento de semillas,
      diferenciando las aproximaciones periódicas utilizadas para
      inicializar la búsqueda de la evolución causal del sistema original.
\item Comparar integradores, pasos, horizontes y tratamientos de memoria
      para determinar qué características de las respuestas persisten y
      cómo cambia el costo de obtenerlas.
\item Caracterizar los destinos y las cuencas mediante trayectorias,
      diagnósticos dinámicos, cortes y sondeos alrededor de todos los
      equilibrios, contrastando la localización estructurada con otras
      estrategias de búsqueda.
\item Relacionar cada resultado con sus parámetros, condiciones iniciales,
      métodos y datos de cálculo, de forma que las figuras y las
      comparaciones puedan reconstruirse a partir de la misma formulación.
\item Estudiar la generación de secuencias y el cifrado reversible de
      imágenes con los sistemas seleccionados, considerando la precisión
      finita, el costo de implementación y la evaluación criptográfica
      del esquema utilizado.
\end{enumerate}

\subsection{Alcance y delimitación}
\label{sec:cap1-alcance}

El núcleo del trabajo comprende sistemas autónomos de dimensión finita
en orden entero y modelos con derivadas de Caputo de orden conmensurable
\(0<q<1\). La formulación de realimentación permite aplicar el análisis
armónico cuando la estructura y el cierre lo admiten. Los métodos
geométricos amplían la comparación a casos en los que esa construcción
directa resulta insuficiente. No se presupone que todas las familias
tengan la misma no linealidad ni que una única semilla pueda localizar
todas sus respuestas.

La selección de casos responde a funciones diferentes dentro del estudio.
Chua permite relacionar las construcciones publicadas con las rutas
fraccionarias propuestas; Kalman--Fitts y PLL permiten examinar
obstrucciones armónicas y alternativas de continuación; MAVPD permite
comparar destinos coexistentes y semillas geométricas. Los casos de
Lorenz generalizado, Rabinovich--Fabrikant y Muñoz--Pacheco aportan
contrastes en otras familias fraccionarias. El interés de incluirlos es
examinar qué elementos del procedimiento se conservan al cambiar el modelo.

La localización se evalúa numéricamente mediante problemas de valor
inicial y protocolos de cuenca especificados. Las afirmaciones sobre
estabilidad, periodicidad y caos se relacionan con el análisis matemático
y los diagnósticos que las sustentan. Los sondeos finitos proporcionan
evidencia numérica de accesibilidad u ocultedad a las escalas estudiadas;
una exclusión matemática de una vecindad completa exige un argumento
adicional sobre esa región. En Caputo, la clase de inicializaciones de
memoria forma parte de la definición del problema de cuenca.

La aplicación digital considera la selección del régimen dinámico, la
generación de la secuencia y la transformación de imágenes como etapas
relacionadas. Su evaluación comprende reproducción, reversibilidad,
estadística y análisis de seguridad del esquema. El comportamiento del
oscilador motiva la construcción y define sus condiciones de operación;
las propiedades criptográficas se estudian sobre el algoritmo que
efectivamente transforma la imagen.

\subsection{Organización de la tesis}
\label{sec:cap1-organizacion}

La tesis se organiza en cinco capítulos. El primero presenta el contexto,
los antecedentes, el problema y los objetivos de la investigación.
El segundo introduce los sistemas de estudio y desarrolla el análisis
de sus ecuaciones, equilibrios y estabilidad. El tercero expone la
formulación teórica y la metodología de localización, continuación e
integración. El cuarto reúne los resultados, las comparaciones entre
métodos y el análisis de las cuencas y de los regímenes observados.
El quinto aborda la implementación digital y la aplicación al cifrado
de imágenes, relacionando sus resultados con el análisis dinámico previo.

\ifHAFOPedagogy
\begin{bloqueestudio}{Lecturas}
Leonov--Kuznetsov desarrolla la inicialización armónica; Danca estudia
Chua con Caputo y Wu utiliza una recurrencia ADM local. La comparación
requiere identificar en cada caso la ecuación de evolución y el dato inicial.
\end{bloqueestudio}

\begin{bloqueestudio}{Ejercicios}
Para el Chua no suave, obtenga el jacobiano en cada equilibrio y relacione su
espectro con el criterio de Matignon. Explique después por qué esa
clasificación local no determina el destino de una semilla armónica alejada
del equilibrio. Repita la comparación para un punto perpetuo.
\end{bloqueestudio}
\fi
